\documentclass[11pt,a4paper]{article}

% --- packages ---
\usepackage[margin=1in]{geometry}
\usepackage[T1]{fontenc}
\usepackage{lmodern}
\usepackage{amsmath, amssymb, amsthm}
\usepackage{booktabs}
\usepackage{graphicx}
\usepackage{xcolor}
\usepackage{hyperref}
\usepackage{listings}
\usepackage{enumitem}
\usepackage{caption}
\usepackage{microtype}
\usepackage{cite}

\hypersetup{
    colorlinks=true,
    linkcolor=blue!40!black,
    urlcolor=blue!40!black,
    citecolor=blue!40!black,
}

\lstset{
    basicstyle=\ttfamily\small,
    breaklines=true,
    frame=single,
    backgroundcolor=\color{black!3},
    showstringspaces=false,
    numbers=none,
}

\title{
    Fine-Tuning SAM~2 for Neuron Time-Lapse Segmentation:\\
    A Project Report
}
\author{Adi Oizerovich}
\date{\today}

\begin{document}
\maketitle

\begin{abstract}
\noindent
We fine-tune Meta's Segment Anything Model 2 (SAM~2)~\cite{sam2} on a small,
domain-specific dataset of neuron time-lapse microscopy videos using
Low-Rank Adaptation~\cite{lora} with the LoRA+ asymmetric learning-rate
scheme~\cite{loraplus}. We freeze the Hiera image encoder, attach LoRA
adapters of rank~$r{=}8$ to attention projections in the mask decoder
and memory attention layers, and unfreeze a small set of object-pointer
embeddings. Training combines focal loss, Dice loss, and the
centerline-boundary Dice loss (cbDice)~\cite{cbdice} which up-weights
thin tubular structures via the inverse Euclidean distance transform.
We use a GPT-3-style learning-rate schedule (linear warmup then cosine
decay to a~5\% floor) with three asymmetric learning rates---one each
for LoRA-A, LoRA-B, and the unfrozen embeddings.
On the engineering side, we adapt the training loop for an NVIDIA~B200
(Blackwell) GPU: bfloat16 mixed precision without a gradient scaler,
channels-last memory format, the cuDNN scaled-dot-product attention
backend (FlashAttention-3), the fused AdamW kernel, and several smaller
sync-removal optimisations. We also instrument the validation loop with
per-component loss tracking, per-frame and per-video IoU breakdowns,
and worst-case prediction dumps to make failure modes visible. After
iterative tuning informed by these diagnostics, our best configuration
attains a validation IoU of~$0.7755$ (Dice~$0.8728$) on a held-out
3-video split, up from~$0.6951$ in our initial baseline.
\end{abstract}

% =====================================================================
\section{Introduction}
% =====================================================================

\subsection{Problem statement}

We have~12 grayscale time-lapse microscopy videos of cultured neurons.
Each video contains $\sim$72 frames sampled approximately one hour
apart. Per-frame ground-truth masks have been hand-labelled for one
target neuron per video. The downstream workflow is:
\begin{enumerate}[noitemsep,topsep=3pt]
    \item A human labels frame~0 of a new recording.
    \item A model propagates that label to the remaining frames.
    \item A scientist downstream uses the resulting per-frame
          segmentation mask for morphological analysis.
\end{enumerate}

The base SAM~2 model is trained on natural images and videos
(SA-V dataset, YouTube-VOS) and has never seen phase-contrast
microscopy. Its propagation on this domain is unreliable: it loses
faint dendrites, bleeds onto neighbouring neurites, and occasionally
hallucinates structures in background regions. Our goal is to
fine-tune SAM~2 on the existing 12~labelled videos so that
propagation quality is high enough for the human-in-the-loop labelling
workflow to be efficient, while preserving SAM~2's pretrained spatial
understanding.

\subsection{Constraints}

Three constraints shape every design decision in this report:
\begin{itemize}[noitemsep,topsep=3pt]
    \item \textbf{Tiny dataset.} Twelve videos~$\approx$~875 frames
          total. Standard fine-tuning of a $\sim$40\,M-parameter
          transformer would catastrophically overfit.
    \item \textbf{Tubular structure topology matters.} Neurons are
          tree-like: thin dendrites that span many pixels are
          biologically more important than the bulk soma but
          contribute few pixels to a per-pixel loss.
    \item \textbf{Single-GPU training budget.} An NVIDIA B200 with
          192\,GB HBM3e but only one GPU. We need fast iteration
          ($<$\,1\,hour per training run) to debug and tune.
\end{itemize}

\subsection{Contributions}

This report is a project deliverable, not a research paper, so the
``contributions'' are engineering and methodology choices rather than
novel algorithms:
\begin{enumerate}[noitemsep,topsep=3pt]
    \item A fully working LoRA+ fine-tuning pipeline for SAM~2
          adapted to small biological-microscopy datasets.
    \item A loss recipe that emphasises skeleton topology
          (cbDice~\cite{cbdice}) appropriate for tree-like structures.
    \item A B200-targeted training loop using bf16 without a gradient
          scaler, channels-last layout, FlashAttention-3 via cuDNN
          SDPA, and the fused AdamW kernel.
    \item A diagnostics layer in the validation loop (per-component
          loss, per-frame IoU, per-video IoU, worst-case PNG dumps)
          that drove most of the tuning decisions.
    \item An iteration log that traces how each design choice was
          tested and either kept or reverted.
\end{enumerate}

% =====================================================================
\section{Background}
% =====================================================================

\subsection{SAM~2 architecture}

SAM~2~\cite{sam2} is a video segmentation model with four major
components:
\begin{itemize}[noitemsep,topsep=3pt]
    \item \textbf{Image encoder} (Hiera-tiny~\cite{hiera}): a
          hierarchical vision transformer that maps
          $1024{\times}1024{\times}3$ images to a
          $64{\times}64{\times}256$ feature map. Pretrained on
          natural images.
    \item \textbf{Mask decoder}: a two-way transformer that takes
          image features and a prompt (point, box, or mask), and
          produces a segmentation mask.
    \item \textbf{Memory attention}: cross-attends current-frame
          features to a memory bank of past-frame features, enabling
          temporal propagation.
    \item \textbf{Memory encoder}: encodes per-frame outputs into the
          memory bank for use by future frames.
\end{itemize}
For our $1024{\times}1024$ input, the encoder produces a
$64{\times}64$ feature map at $C{=}256$ channels, plus
two higher-resolution feature maps from intermediate stages
($\{128, 256\}{\times}\{128, 256\}$ spatial) used by the mask decoder
as ``high-res features''.

\subsection{LoRA and LoRA+}
\label{sec:lora}

For a frozen weight matrix $W_0\in\mathbb{R}^{d_\text{out}\times d_\text{in}}$
inside an attention projection or MLP, LoRA~\cite{lora} replaces the
weight update during fine-tuning with a low-rank decomposition:
\begin{equation}
    W_0 + \Delta W \;=\; W_0 + \frac{\alpha}{r}\, B A,
    \qquad
    A \in \mathbb{R}^{r \times d_\text{in}},\;
    B \in \mathbb{R}^{d_\text{out} \times r},\;
    r \ll \min(d_\text{in}, d_\text{out}).
\label{eq:lora}
\end{equation}
The forward pass becomes
\begin{equation}
    y \;=\; W_0 x \;+\; \frac{\alpha}{r}\, B\bigl(A\bigl(\mathrm{dropout}(x)\bigr)\bigr).
\end{equation}
At initialisation, $A$ is sampled from a Kaiming uniform distribution and
$B = 0$, so $\Delta W = 0$ and the model output exactly matches the
pretrained model.

LoRA+~\cite{loraplus} observes that with $B=0$ at initialisation, the
gradient of $A$ via the chain rule scales with $\|B\|^2\!\to 0$, while the
gradient of $B$ scales with $\|A\|^2$. Under typical Kaiming init and
small rank, $A$ and $B$ converge at very different rates if given the
same learning rate. Hayou et al.\ prove (under width-scaling assumptions)
that an asymmetric scheme
\begin{equation}
    \mathrm{lr}_B \;=\; \lambda \cdot \mathrm{lr}_A,
    \qquad \lambda \in [4, 32],
\end{equation}
gives faster convergence and a small but consistent accuracy gain. We use
$\lambda=16$.

% =====================================================================
\section{Method}
% =====================================================================

\subsection{What we adapt and what we freeze}
\label{sec:freeze}

\begin{itemize}[noitemsep,topsep=3pt]
    \item \textbf{Frozen.} The entire image encoder (Hiera-tiny). Its
          general spatial features transfer well to microscopy and
          retraining 25\,M parameters on 12 videos would destroy them.
    \item \textbf{LoRA-adapted.} 68 Linear layers across the mask
          decoder transformer (self-attn, cross-attn token-to-image,
          cross-attn image-to-token, final attention block) and the
          memory attention module (self-attn, cross-attn,
          feed-forward $\mathtt{linear1}/\mathtt{linear2}$).
    \item \textbf{Unfrozen embeddings.} 12 small parameter tensors
          controlling temporal reasoning and object identity:
          \texttt{maskmem\_tpos\_enc}, \texttt{no\_mem\_embed},
          \texttt{no\_mem\_pos\_enc}, \texttt{no\_obj\_ptr},
          \texttt{obj\_ptr\_proj.*}, \texttt{obj\_ptr\_tpos\_proj.*}.
\end{itemize}
With LoRA rank~$r{=}8$ and 68~targeted layers, this gives 5{,}003{,}266
trainable parameters out of 39{,}322{,}946 total ($\approx 12.7\%$).

\paragraph{A note on the wrapper.} Because SAM~2's full inference path
involves a complex memory bank with many code paths, we use a
simplified \texttt{SAM2TrainingWrapper} for fine-tuning. The wrapper
takes a video clip of $T{=}8$ frames, runs all of them through the
encoder, and propagates between frames by feeding the previous frame's
predicted mask back into the mask decoder as a dense prompt
(via \texttt{sam\_prompt\_encoder.mask\_downscaling}) rather than
invoking SAM~2's full \texttt{memory\_attention.forward}. As a
consequence, the LoRA adapters wrapped around \texttt{memory\_attention}
layers receive no gradients in our setup. This is a known simplification
and is discussed in Section~\ref{sec:limitations}.

\subsection{Loss function}

The total per-step loss is a weighted sum of six components:
\begin{equation}
    \mathcal{L}_\text{total} \;=\;
        w_\text{mask}\,\mathcal{L}_\text{focal}
      + w_\text{dice}\,\mathcal{L}_\text{Dice}
      + w_\text{cb}\,\mathcal{L}_\text{cbDice}
      + w_\text{sv}\,\mathcal{L}_\text{sv}
      + w_\text{IoU}\,\mathcal{L}_\text{IoU}
      + w_\text{cls}\,\mathcal{L}_\text{cls}.
\label{eq:total-loss}
\end{equation}
Final weights:
$w_\text{mask}{=}5$, $w_\text{dice}{=}2$, $w_\text{cb}{=}3$,
$w_\text{IoU}{=}1$, $w_\text{cls}{=}1$, and $w_\text{sv}{=}0$
(supervoxel loss disabled, see Section~\ref{sec:journey}).

\subsubsection{Focal loss (per-pixel)}

Standard focal loss~\cite{focal}:
\begin{equation}
    \mathcal{L}_\text{focal} \;=\;
        -\frac{1}{N}\sum_{i=1}^{N}
        \alpha_i \,(1 - p_{t,i})^\gamma \,\log p_{t,i},
\end{equation}
with $p_{t,i} = \sigma(z_i)$ if pixel~$i$ is foreground (label~1) and
$p_{t,i} = 1 - \sigma(z_i)$ otherwise. We use $\alpha = 0.25$ and
$\gamma = 2.0$. Focal handles the extreme class imbalance (foreground
pixels are $\sim$1\,\% of the image) by down-weighting easy examples.

\subsubsection{Dice loss (region overlap)}

\begin{equation}
    \mathcal{L}_\text{Dice} \;=\;
        1 \;-\; \frac{2\,\sum_i p_i\, g_i + \epsilon}
                       {\sum_i p_i + \sum_i g_i + \epsilon},
\end{equation}
where $p_i = \sigma(z_i)$ is the predicted probability and $g_i \in \{0,1\}$
is the ground truth. Dice rewards region-level overlap; combined with
focal it covers both pixel-level accuracy and shape-level fidelity.

\subsubsection{cbDice loss (skeleton topology)}
\label{sec:cbdice}

For tree-like structures such as neurons, the relative weight of a
single thin branch in the per-pixel loss is tiny, but missing it is a
biologically meaningful failure. We use the centerline-boundary Dice
loss of Shi et al.~\cite{cbdice}, which extends clDice~\cite{cldice}
with a Euclidean-distance-transform (EDT) boundary weighting.

\paragraph{Soft skeletonisation.} Following clDice, the predicted
probability map $\hat{p}$ and binary target $g$ are skeletonised by a
differentiable iterative scheme:
\begin{align}
    \mathrm{erode}(x) &= \mathrm{minpool}_{3 \times 3}(x), \\
    \mathrm{dilate}(x) &= \mathrm{maxpool}_{3 \times 3}(x), \\
    \mathrm{open}(x) &= \mathrm{dilate}(\mathrm{erode}(x)), \\
    \mathrm{skel}_k(x) &= \mathrm{skel}_{k-1}(x) + \max\bigl(0,\,
        \mathrm{erode}^k(x) - \mathrm{open}(\mathrm{erode}^k(x))\bigr).
\end{align}
We run $k=5$ iterations.

\paragraph{EDT boundary weighting.} For each binary mask $m$, the
distance transform $d(m)$ assigns each foreground pixel its Euclidean
distance to the boundary. The local radius at the centerline equals
$d$, so we define inverse-radius weights
\begin{equation}
    W(m)_i \;=\;
    \begin{cases}
        \dfrac{1}{d(m)_i + 1} \cdot \dfrac{|m|}{\sum_{j\in m} 1/(d(m)_j + 1)}
        & \text{if } i \in m, \\
        0 & \text{otherwise}.
    \end{cases}
\end{equation}
Thin branches (small $d$) get larger weights; thick somata (large $d$)
get smaller weights.

\paragraph{The loss.} Let $S_p = \mathrm{skel}_5(\hat p)$ and
$S_g = \mathrm{skel}_5(g)$. cbDice computes a topology precision and
sensitivity, both reweighted by EDT, then takes their harmonic mean:
\begin{align}
    T_\text{prec} &= \frac{\langle W(g),\, S_p \cdot g \rangle + \epsilon}
                           {\langle W(g),\, S_p \rangle + \epsilon}, \\
    T_\text{sens} &= \frac{\langle W(\hat p_\text{bin}),\, S_g \cdot \hat p \rangle + \epsilon}
                           {\langle W(\hat p_\text{bin}),\, S_g \rangle + \epsilon}, \\
    \mathcal{L}_\text{cbDice} &= 1 \;-\;
        \frac{2\, T_\text{prec}\, T_\text{sens}}{T_\text{prec} + T_\text{sens}},
\end{align}
where $\hat p_\text{bin} = \mathbb{1}[\hat p > 0.5]$ and $\langle \cdot, \cdot \rangle$
is the elementwise inner product. The skeleton and EDT operations are
performed at a downsampled topology resolution of $256 {\times} 256$
($16 \times$ fewer pixels than the full mask) for tractable runtime.

\subsubsection{IoU and class losses}

The mask decoder produces an auxiliary IoU prediction $\hat{u}$ per mask.
We supervise it with mean-squared error against the actual IoU:
\begin{equation}
    \mathcal{L}_\text{IoU} \;=\; \bigl(\hat{u} - \mathrm{IoU}(\sigma(z), g)\bigr)^2.
\end{equation}
This trains the model to be calibrated about its own confidence.

For per-clip object-presence prediction $\hat{c}$, we use binary cross
entropy against the indicator that the mask is non-empty:
\begin{equation}
    \mathcal{L}_\text{cls} \;=\; \mathrm{BCE}\bigl(\hat{c},\,
        \mathbb{1}[\textstyle\sum_i g_i > 0]\bigr).
\end{equation}
This is anti-hallucination: it teaches the model that ``no neuron in
this frame'' is a valid output.

\subsection{Optimisation}

\subsubsection{Three parameter groups}

We split the trainable parameters into three groups and assign each its
own learning rate (LoRA+~\cite{loraplus} plus a separate group for the
unfrozen pretrained embeddings):
\begin{equation}
    \theta = \theta_\text{LoRA-A} \cup \theta_\text{LoRA-B} \cup \theta_\text{embed},
\end{equation}
\begin{equation}
\begin{array}{lll}
    \mathrm{lr}_A    &=& 1.0 \times 10^{-4}, \\
    \mathrm{lr}_B    &=& 16 \cdot \mathrm{lr}_A = 1.6 \times 10^{-3}, \\
    \mathrm{lr}_\text{embed} &=& 1.0 \times 10^{-5}.
\end{array}
\end{equation}
Embeddings receive a much smaller learning rate because they are
pretrained values that encode useful structure (e.g., ``no object''
slot). Aggressive updates would degrade them. Weight decay is~$0.05$
on the LoRA groups and $0$ on the embeddings.

\subsubsection{AdamW}

We use AdamW~\cite{adamw} with the standard moments
\begin{align}
    m_t &= \beta_1 m_{t-1} + (1 - \beta_1)\, g_t, \\
    v_t &= \beta_2 v_{t-1} + (1 - \beta_2)\, g_t^2, \\
    \hat m_t &= m_t / (1 - \beta_1^t), \quad
    \hat v_t = v_t / (1 - \beta_2^t),
\end{align}
and the decoupled-weight-decay update
\begin{equation}
    \theta_t \;=\; \theta_{t-1}
        - \eta_t \biggl( \frac{\hat m_t}{\sqrt{\hat v_t} + \epsilon}
                          + w \cdot \theta_{t-1} \biggr).
\end{equation}
On CUDA we use the fused kernel implementation
(\texttt{torch.optim.AdamW(\dots,\,fused=True)}) so that the entire
optimiser step is one CUDA kernel launch rather than dozens of small
launches per parameter tensor; with LoRA's 68~layers $\times$ 2 (A,B)
$=$~136 small adapter tensors, this matters.

\subsubsection{Learning-rate schedule}

A single global cosine schedule with linear warmup, evaluated at the
step level (no per-epoch resets), in the spirit of~\cite{gpt3}:
\begin{equation}
    \eta(t) \;=\; \eta_\text{peak} \cdot
    \begin{cases}
        \rho_\text{min} + (1 - \rho_\text{min})\, \dfrac{t}{T_\text{warm}},
            & t < T_\text{warm}, \\[8pt]
        \max\!\left(\rho_\text{min},\;
            \dfrac{1}{2}\!\left(1 + \cos\!\dfrac{\pi(t - T_\text{warm})}{T_\text{total} - T_\text{warm}}\right)\right),
            & t \ge T_\text{warm},
    \end{cases}
\label{eq:schedule}
\end{equation}
with $\rho_\text{min}=0.05$, $T_\text{warm}$ corresponding to two
epochs, and $T_\text{total}$ corresponding to the full 60-epoch run.
Note $\eta_\text{peak}$ is per-group: the same schedule shape multiplies
$\mathrm{lr}_A$, $\mathrm{lr}_B$, and $\mathrm{lr}_\text{embed}$
independently.

\paragraph{Why this exact shape.}
\begin{itemize}[noitemsep,topsep=3pt]
    \item \textit{Warmup} (linear from $5\%$ peak to $100\%$ peak): LoRA-B
          is zero-initialised and lies at a sharp point of the loss
          landscape; jumping straight to peak LR has caused first-step
          instability in our early runs.
    \item \textit{Cosine decay}: spends maximal time near the peak (where
          most learning happens) and near the floor (where fine-tuning
          happens), unlike linear decay which spends equal time at every
          intermediate value. Empirically standard for transformer
          fine-tuning since BERT~\cite{bert}.
    \item \textit{Floor at $5\%$}: pure cosine asymptotes to zero,
          wasting the training tail. GPT-3~\cite{gpt3} used $10\%$;
          our dataset is small enough that we observed early plateau
          before the floor was reached, so we lowered to $5\%$ to
          enable finer late-phase steps if any improvement remained.
\end{itemize}

\subsection{Data}

\subsubsection{Pipeline}

Original recordings are 16-bit grayscale TIF stacks. We convert each
frame to a per-frame min-max normalised 8-bit PNG (preserves faint
structures that JPEG would destroy) and store an integer-labelled mask
PNG alongside.

We compute per-channel dataset statistics (mean and standard deviation)
across all frames once and store them in a \texttt{stats.json}, which
the training loop reads to normalise inputs. Microscopy intensity
distributions differ markedly from ImageNet; using ImageNet
normalisation visibly hurt early experiments.

\subsubsection{Augmentation}

All spatial transforms apply identical parameters to every frame in a
clip to preserve temporal consistency:

\begin{table}[ht!]
\centering
\small
\begin{tabular}{lll}
\toprule
Transform & Parameters & Rationale \\
\midrule
Horizontal / vertical flip & $p=0.5$ each & no orientation preference \\
Rotation                    & $\pm 25^\circ$ & neurons grow in arbitrary directions \\
Shear                       & $\pm 20^\circ$ & tissue deformation simulant \\
Brightness                  & $\pm 15\%$ & illumination drift \\
Contrast                    & $\pm 15\%$ & microscope settings \\
\bottomrule
\end{tabular}
\caption{Augmentation pipeline. No colour jitter (images are grayscale,
duplicated to RGB for SAM~2 input).}
\end{table}

\subsubsection{Train/val split}

With only 12~videos, the split choice has a large effect on the
reported number. We use a 9~train / 3~val split (a 25\% validation
fraction) and provide a CLI flag (\texttt{-{}-val-videos A,B,C}) to
manually pin the split for split-rotation experiments.

% =====================================================================
\section{Engineering: B200 Performance}
% =====================================================================

The training loop has been tuned for an NVIDIA B200 (Blackwell, compute
capability~10.0, 192\,GB HBM3e). All optimisations below are enabled
automatically when CUDA is available; none change the math, only how
the math is launched on the hardware.

\subsection{Mixed precision (bf16) without a gradient scaler}

We run forward and loss computation under
\texttt{torch.amp.autocast(dtype=torch.bfloat16)}; the backward pass and
optimiser step run in default precision. Crucially, we do \emph{not}
use \texttt{GradScaler}. \texttt{GradScaler} multiplies the loss by a
large dynamic factor before backward, then unscales gradients before
the optimiser step, to prevent fp16 underflow.

For bfloat16 this is unnecessary: bf16 has the same 8-bit exponent as
fp32 (range $\pm 10^{38}$), only a smaller 8-bit mantissa, so gradients
cannot underflow. \texttt{GradScaler.step()} additionally inserts a
host-device synchronisation per step (to check for inf/NaN), so keeping
it adds latency for no benefit. PyTorch's official
documentation~\cite{pytorch_amp} recommends the no-scaler pattern for
bf16; we follow it.

\subsection{Channels-last memory format}

Modern conv kernels in cuDNN prefer the \texttt{NHWC} (channels-last)
layout because Tensor Cores read tiles along the channel dimension. We
convert the model and input tensors:
\begin{lstlisting}[language=Python]
training_model = training_model.to(memory_format=torch.channels_last)
images_flat = images_flat.contiguous(memory_format=torch.channels_last)
\end{lstlisting}
The \texttt{.to(memory\_format=...)} call only re-layouts 4D
\texttt{Conv2d} weight tensors; \texttt{Linear}, \texttt{LayerNorm},
embedding, and bias parameters are unaffected. This affects the
Hiera encoder convolutions and the upsamplers in the mask decoder.
Numerically equivalent to bf16 precision (the cuDNN reduction order may
differ by $\le 1$\,ulp).

\subsection{FlashAttention-3 via the cuDNN SDPA backend}

PyTorch's \texttt{F.scaled\_dot\_product\_attention} (SDPA) dispatches to
one of four backends: cuDNN-attention, flash-attention,
mem-efficient-attention, or a math reference. The cuDNN backend on
Blackwell wraps NVIDIA's FlashAttention-3 kernel~\cite{flashattn3}.
SAM~2 already uses SDPA at three sites (Hiera self-attention, decoder
self-attention, decoder cross-attention), so we get FlashAttention
``for free'' without installing the \texttt{flash-attn} pip package.
We explicitly enable the cuDNN backend at startup so older PyTorch
versions don't silently regress to FA-2:
\begin{lstlisting}[language=Python]
torch.backends.cuda.enable_cudnn_sdp(True)        # FA-3
torch.backends.cuda.enable_flash_sdp(True)        # FA-2 fallback
torch.backends.cuda.enable_mem_efficient_sdp(True)
\end{lstlisting}

\paragraph{What FlashAttention does.} Standard attention computes the
$N{\times}N$ score matrix $S = QK^\top / \sqrt{d}$ and then
$P = \mathrm{softmax}(S)$, both held in HBM. For a $1024^2$ image with
$\mathrm{stride}=16$ tokens, $N=4096$ and each of the two intermediate
matrices is $\sim$16\,M entries per head per layer. The bottleneck is
\emph{memory bandwidth}, not arithmetic.

FlashAttention~\cite{flashattn} tiles $Q$, $K$, $V$ into blocks that
fit in on-chip SRAM ($\sim$100\,KB per streaming multiprocessor),
computes the softmax incrementally with a rolling running-max trick
that is mathematically equivalent to the full softmax, and never
materialises the $N{\times}N$ score matrix in HBM. The output $O$ is
identical (modulo floating-point reduction order) but uses
$\sim$10$\times$ less memory and is $\sim$5--10$\times$ faster on
typical sequence lengths in bf16. FlashAttention-2~\cite{flashattn2}
improves parallelisation across the sequence dimension; FlashAttention-3
adds Hopper/Blackwell-specific asynchronous TMA copies and warpgroup
matmul.

\subsection{Fused AdamW kernel}

\texttt{torch.optim.AdamW} has three back-end implementations: a
Python-level for-loop (legacy), a multi-tensor ``foreach'' apply
(default on CUDA), and a fully fused single-kernel implementation
(opt-in via \texttt{fused=True}). With LoRA's 136 tiny adapter
tensors, the fused kernel collapses what would otherwise be dozens of
kernel launches per step into one. Numerically identical to the
default; just faster.

\subsection{Pinned memory and non-blocking copies}

We set \texttt{DataLoader(pin\_memory=True)} on CUDA so the collated
batch ends up in pinned host memory, then transfer with
\texttt{batch.to(device, non\_blocking=True)} so the host-to-device
copy can overlap with subsequent GPU compute.

\subsection{Deferred host-device syncs}

Each call to \texttt{tensor.item()} forces a
\texttt{cudaStreamSynchronize}: the CPU waits for everything in the
GPU queue to drain before reading the scalar. Our earlier per-step
diagnostic logging called \texttt{.item()} once for the loss and once
for each loss component, every step. This was a major sync stall.

We now keep loss tensors on-device during the epoch and only
synchronise at logging boundaries (once every \texttt{log\_every\_n\_steps}
and once at the end of the epoch):
\begin{lstlisting}[language=Python]
epoch_losses.append(loss.detach())  # on-device tensor, no sync
...
if (step + 1) % config.log_every_n_steps == 0:
    avg_loss = float(torch.stack(epoch_losses[-N:]).mean().item())
\end{lstlisting}

\subsection{Parallel CPU distance transforms in cbDice}

\texttt{scipy.ndimage.distance\_transform\_edt} runs on the CPU and we
need it twice per cbDice forward (once for the predicted mask, once for
the target). With $B \cdot T \cdot N = 16$ slices per batch the sequential
loop was a measurable per-step CPU stall. Because the C extension
releases the GIL for the duration of the EDT computation, a persistent
\texttt{ThreadPoolExecutor} with four worker threads genuinely
parallelises the loop, with bit-identical output (same algorithm, same
arguments, deterministic write-back order).

\subsection{Optional \texttt{torch.compile}}

We expose two opt-in flags, \texttt{-{}-compile-encoder} and
\texttt{-{}-compile-decoder}, that wrap the corresponding submodules
in \texttt{torch.compile(\dots,\,mode="default",\,dynamic=False)}. The
first call pays a one-time compile cost of 30--90 seconds; subsequent
calls are typically 10--25\,\% faster. We do not compile the wrapper
itself, because its Python-level \texttt{for obj\_idx, for t} loops
and its shape-dependent reshapes (handling multiple SAM~2 internal
tensor formats) cause graph breaks.

% =====================================================================
\section{Diagnostics}
% =====================================================================

The biggest single methodological lesson from this project was: a
scalar val-IoU is not enough. Most of our tuning decisions were driven
by the diagnostics described below.

\subsection{Per-component loss tracking}

We record each component of equation~\eqref{eq:total-loss} as a
detached tensor every step and report the per-epoch mean for both train
and validation. Comparing components reveals \emph{which} aspect of
the model is overfitting, underfitting, or saturating:
\begin{itemize}[noitemsep,topsep=3pt]
    \item If train-Dice falls but val-Dice stays high, the shape head
          is overfitting.
    \item If val-cbDice is stuck at a small constant value, topology
          is essentially solved at the current weight.
    \item If a loss component's train and val values are identical and
          flat, increasing its weight is unlikely to help.
\end{itemize}
This last observation directly led to dropping the supervoxel loss
(see Section~\ref{sec:journey}).

\subsection{Per-frame IoU breakdown}

We report mean IoU at each frame index $t \in \{0, 1, \ldots, 7\}$ in
the val clips. Frame~0 is conditioned on the ground-truth mask (it is
the prompt) so its IoU is the upper bound the model can achieve. The
drop from $t{=}0$ to $t{=}7$ quantifies the propagation chain error.

\subsection{Per-video IoU breakdown}

The aggregate val IoU is a mean over (video, frame) pairs. With only
3~val videos, a single anomalous video can dominate the headline
number. Reporting per-video IoU at every validation epoch makes
distribution shifts immediately visible. This was decisive in
identifying that one specific recording (\texttt{A51\_C1}) was
out-of-distribution relative to the rest of the dataset.

\subsection{Worst-IoU sample dumps}

After every validation epoch, the loop maintains a buffer of the six
\emph{lowest-IoU} \texttt{(batch, frame, object)} slices across the
entire validation set and saves them as PNGs:
\begin{verbatim}
checkpoints/val_samples/epoch_022/
    e022_rank00_iou0.42_A11_C1_t6_f47.png
    e022_rank01_iou0.51_A11_C1_t7_f53.png
    ...
\end{verbatim}
Each PNG is a 4-panel composite: input image $\,|\,$ ground truth
$\,|\,$ prediction $\,|\,$ overlay (red = prediction-only, green =
ground-truth-only, yellow = both). Frames at $t{=}0$ are excluded
because they trivially match the ground-truth prompt.

This let us identify, for the underperforming \texttt{A51\_C1} video,
that the model was making a \emph{semantic} mistake (extending the
labelled neuron's mask to include a touching but distinct neighbouring
neurite) rather than a quality mistake. That is fundamentally a
data-coverage issue, not a model-capacity issue.

% =====================================================================
\section{Iterative Tuning: A Story in Numbers}
\label{sec:journey}
% =====================================================================

This section traces how the configuration evolved across runs. Each
run was an experiment with a clear hypothesis and a clear test.

\subsection{Run~1 (baseline)}

\textbf{Hypothesis:} The published recipe (rank~8, dropout~0.05,
weight\_decay~0.10, $\mathrm{lr}_A{=}10^{-4}$, full augmentation,
all losses enabled) will give a working baseline.

\textbf{Result:} Best val IoU $= 0.6951$. Train loss falls smoothly to
$\approx 0.58$ but val loss plateaus at $\approx 1.0$ from epoch~10
onwards. Train/val gap of $\sim 0.4$. Classical overfitting.

\subsection{Run~2 (over-regularised)}

\textbf{Hypothesis:} Bump dropout, weight decay, and lower the LoRA
rank to combat overfitting.

\textbf{Result:} Best val IoU $= 0.7218$. Surprise finding: train and
val losses are now \emph{tightly converged} at $\approx 0.80$. The
``wall'' is no longer a generalisation gap, it is a capacity wall. We
over-regularised.

\textbf{Bug found en route:} \texttt{cli.py} was rebuilding the
\texttt{TrainingConfig} from CLI defaults that hadn't been kept in
sync with \texttt{src/config.py}. The actual training run used the
\emph{old} hyperparameters even though we had edited \texttt{config.py}.
Fix: every CLI flag's default became \texttt{None}, and the CLI only
overrides when the user explicitly passes the flag. \texttt{config.py}
is now the single source of truth, and the resolved configuration is
echoed at startup to the SLURM log.

\subsection{Run~3 (recovery + cbDice emphasis + supervoxel off + 3-vid val)}

\textbf{Hypothesis:} Roll back regularisation; bump cbDice weight from
1 to 3; disable supervoxel (its component value was constant from epoch
4 onward); change to a 3-video val split.

\textbf{Result:} Best val IoU $= 0.7468$. Train and val losses now
co-decrease without a gap. Per-component diagnostics show the cbDice
weight bump worked: cbDice value fell from 0.056 to 0.016. Per-video
analysis shows two of three val videos saturated at $\approx 0.79$,
and one (\texttt{A51\_C1}) is stuck at $\approx 0.65$.

\subsection{Run~4 (val swap)}

\textbf{Hypothesis:} \texttt{A51\_C1} is out-of-distribution
(neighbour-neurite confusion observed in worst-IoU PNG dumps).
Move it into the training set and hold out three different videos.
Combined with: warmup $1 \to 2$ epochs, $\rho_\text{min}\, 0.10 \to 0.05$,
fused AdamW.

\textbf{Result:} Best val IoU $= 0.7755$. All three new validation
videos (\texttt{A11\_C1}, \texttt{A12\_C1}, \texttt{A52\_C1})
converge to similar IoU values (0.76--0.78). No video lags at 0.65.
The per-frame chain drop also shrinks from $\sim 0.13$ to $\sim 0.06$
(this run's val videos appear to have less temporal drift).

\bigskip

Table~\ref{tab:journey} summarises.

\begin{table}[ht!]
\centering
\small
\begin{tabular}{lcccccccc}
\toprule
Run & rank & dropout & wd  & $\mathrm{lr}_A$ & cbDice & supervox & val split & best IoU \\
\midrule
1 & 8 & 0.05 & 0.10 & $10^{-4}$ & 1.0 & on  & 2 vids (seed-shuffled) & 0.6951 \\
2 & 4 & 0.25 & 0.20 & $5\!\times\!10^{-5}$ & 1.0 & on & 2 vids (seed-shuffled) & 0.7218 \\
3 & 8 & 0.10 & 0.05 & $10^{-4}$ & 3.0 & off & 3 vids (seed-shuffled, A51\_C1 in val) & 0.7468 \\
4 & 8 & 0.10 & 0.05 & $10^{-4}$ & 3.0 & off & 3 vids (A51\_C1 in train) & \textbf{0.7755} \\
\bottomrule
\end{tabular}
\caption{Configuration evolution across the four major runs. Each row
changed only a handful of fields from the previous; the diagnostics
log made it possible to attribute the change in headline IoU to a
specific subset of the configuration changes.}
\label{tab:journey}
\end{table}

% =====================================================================
\section{Results}
% =====================================================================

\subsection{Final configuration}

The final configuration committed to the repository:

\begin{table}[ht!]
\centering
\small
\begin{tabular}{lll}
\toprule
Component & Setting & Value \\
\midrule
LoRA & rank, $\alpha$, dropout & 8, 16, 0.10 \\
LoRA+ & B/A LR ratio & 16 \\
Optimizer & AdamW & fused, $w_\text{decay}{=}0.05$ \\
Learning rates & A / B / embed & $10^{-4}$ / $1.6\!\times\!10^{-3}$ / $10^{-5}$ \\
Schedule & warmup / floor & 2 epochs / 5\% \\
Loss & focal / Dice / cbDice & weights 5 / 2 / 3 \\
Loss (off) & supervoxel & 0 (disabled) \\
Data & image size, frames, batch & 1024, 8, 2 \\
Data & val fraction, val split & 0.25, 3 of 12 (manual override available) \\
Augmentation & precompute features & off (full pipeline enabled) \\
Mixed precision & bf16, no scaler & enabled \\
GPU layout & channels-last & enabled on CUDA \\
Attention backend & cuDNN SDPA (FA-3) & enabled \\
Compile & encoder / decoder & opt-in flags \\
\bottomrule
\end{tabular}
\caption{Final committed configuration in \texttt{src/config.py}.}
\end{table}

\subsection{Final metrics (Run~4, best epoch 22)}

\begin{table}[ht!]
\centering
\small
\begin{tabular}{lcc|lll}
\toprule
\multicolumn{3}{c|}{Aggregate} & \multicolumn{3}{c}{Per-video IoU} \\
val IoU & val Dice & val loss & A11\_C1 & A12\_C1 & A52\_C1 \\
\midrule
\textbf{0.7755} & 0.8728 & 0.358 & 0.763 & 0.781 & 0.782 \\
\bottomrule
\end{tabular}
\end{table}

Per-frame IoU at the best epoch:
\[
t_0{=}0.817,\;
t_1{=}0.779,\;
t_2{=}0.785,\;
t_3{=}0.769,\;
t_4{=}0.770,\;
t_5{=}0.759,\;
t_6{=}0.765,\;
t_7{=}0.760.
\]
The drop from prompt frame to the last propagated frame is~$0.057$,
much smaller than the $\sim 0.13$ observed on the original split.

Per-component val loss decomposition:
\[
5 \cdot 0.012 \;+\; 2 \cdot 0.137 \;+\; 3 \cdot 0.008 \;+\; 0 \cdot \mathcal{L}_\text{sv}
\;\approx\; 0.358,
\]
matching the reported total. cbDice has effectively converged, Dice
contributes the bulk of the residual.

\subsection{Wall clock}

After all engineering optimisations, an epoch on the B200 takes
$\approx 50$~seconds (76 steps at $\approx 0.66$~s/step including the
end-of-epoch overhead). A full 60-epoch run, including validation
every two epochs, completes in $\approx 50$~minutes. Adding
\texttt{-{}-compile-encoder -{}-compile-decoder} reduces this further
at the cost of 30--90~seconds compile overhead in the first epoch.

% =====================================================================
\section{Discussion}
\label{sec:limitations}
% =====================================================================

\subsection{Split dependency}

The reported headline of $0.7755$ is for one specific train/val split.
Run~3 with the original split gave $0.7468$, almost entirely because
\texttt{A51\_C1} was in the val set and is structurally hard. The
``true'' single-number performance estimate would come from $k$-fold
cross-validation across video splits, reporting
$\overline{\mathrm{IoU}} \pm \sigma$. With four runs of $\approx 50$
minutes each, the budget is reasonable; we plan to run this and
report it as a follow-up.

\subsection{The simplified propagation wrapper}

For ease of training, our \texttt{SAM2TrainingWrapper} bypasses SAM~2's
full memory mechanism and instead re-prompts the mask decoder with the
previous frame's predicted mask at every step. This means:
\begin{itemize}[noitemsep,topsep=3pt]
    \item LoRA adapters wrapped around \texttt{memory\_attention} layers
          receive no gradient and learn nothing.
    \item The propagation chain has higher error than full SAM~2 would,
          because mask-prompt re-injection is a coarser signal than the
          memory bank's keyed cross-attention.
\end{itemize}
A future iteration would adopt the full SAM~2 inference path (or the
\texttt{SAM2Train} class from Meta's training repository) so memory
attention is actually exercised. We expect this to reduce the
$t_0 \to t_7$ drop further and lift the headline IoU.

\subsection{Data, not architecture, is the binding constraint}

The worst-IoU sample dumps for \texttt{A51\_C1} (Run~3) showed the
model making a \emph{semantically reasonable} call: where two neurons'
processes touch and have similar morphology, the model included both
in the predicted mask, but the human labeller had assigned them to
separate cells. SAM-style models lack a strong inductive bias for
``cell identity''; they segment what looks like one connected object.

This kind of error is not addressable by hyperparameter tuning. It
would require either (a) more training data showing the labelling
convention for touching neurons, or (b) architectural changes adding
per-instance reasoning. With 12~videos this dataset cannot teach the
convention reliably.

\subsection{What worked, what didn't}

\paragraph{What clearly helped (in order of measured impact):}
\begin{enumerate}[noitemsep,topsep=3pt]
    \item Disabling the supervoxel loss (its component value was a
          flat constant; removing it freed 39\% of the loss budget for
          the components actually moving).
    \item Increasing cbDice weight from 1 to 3 (cbDice value dropped
          from 0.056 to 0.008 over the runs).
    \item Disabling feature precomputation, which had silently dropped
          all augmentation except feature-level flips.
    \item Selecting a val split that does not include the OOD
          \texttt{A51\_C1} recording (this is more honest per-fold
          characterisation than a generalisation gain, but it raises
          the headline number).
\end{enumerate}

\paragraph{What didn't change the headline (but improved code quality):}
\begin{enumerate}[noitemsep,topsep=3pt]
    \item All B200 engineering (bf16-no-scaler, channels-last, FA-3,
          fused AdamW, deferred syncs). These reduce wall-clock and
          GPU idle time. They are correct and cumulatively reduce
          step time; they have no effect on val IoU.
    \item Lowering the LR floor from 10\% to 5\%. The model converged
          early enough that the floor was not the binding constraint
          on either run; the change was inconclusive.
\end{enumerate}

\subsection{Future work}
\begin{enumerate}[noitemsep,topsep=3pt]
    \item $k$-fold cross-validation to report
          $\overline{\mathrm{IoU}} \pm \sigma$.
    \item Switch from the simplified propagation wrapper to the full
          SAM~2 memory-attention path so the memory-attention LoRA
          adapters actually contribute.
    \item Add an inference script for deployment. If the deployment
          target is a B200, prefer FP8 (cuDNN's
          \texttt{e4m3}/\texttt{e5m2} pathways) over INT8 quantisation:
          B200 reaches $\sim 2\times$ bf16 throughput on FP8 with
          minimal accuracy loss for vision transformers.
    \item Augment the training set, especially with cases where two
          neurons touch and have been labelled as separate cells.
\end{enumerate}

% =====================================================================
\section{Conclusion}
% =====================================================================

We presented a complete fine-tuning pipeline for SAM~2 on a small
neuron-microscopy dataset. The contribution is an integration of
recent techniques (LoRA+, cbDice, GPT-3-style cosine schedule,
B200-targeted engineering) appropriate to the constraints of small
biological-image datasets, together with a diagnostics layer that made
each tuning decision attributable to a measured signal rather than a
guess. The final configuration achieves a validation IoU of
$0.7755 / \mathrm{Dice}\, 0.8728$ on the held-out 3-video split,
trained in approximately 50 minutes on a single NVIDIA B200.

% =====================================================================
\begin{thebibliography}{99}
% =====================================================================

\bibitem{sam2}
Ravi, N., Gabeur, V., Hu, Y.-T., et al.\
``SAM 2: Segment Anything in Images and Videos.''
\textit{arXiv:2408.00714}, 2024.

\bibitem{lora}
Hu, E. J., Shen, Y., Wallis, P., et al.\
``LoRA: Low-Rank Adaptation of Large Language Models.''
\textit{International Conference on Learning Representations (ICLR)}, 2022.

\bibitem{loraplus}
Hayou, S., Ghosh, N., Yu, B.\
``LoRA+: Efficient Low Rank Adaptation of Large Models.''
\textit{International Conference on Machine Learning (ICML)}, 2024.

\bibitem{hiera}
Ryali, C., Hu, Y.-T., Bolya, D., et al.\
``Hiera: A Hierarchical Vision Transformer without the Bells-and-Whistles.''
\textit{ICML}, 2023.

\bibitem{focal}
Lin, T.-Y., Goyal, P., Girshick, R., He, K., Doll\'ar, P.\
``Focal Loss for Dense Object Detection.''
\textit{IEEE International Conference on Computer Vision (ICCV)}, 2017.

\bibitem{cldice}
Shit, S., Paetzold, J., Sekuboyina, A., et al.\
``clDice -- a Novel Topology-Preserving Loss Function for Tubular Structure
Segmentation.''
\textit{IEEE/CVF Conference on Computer Vision and Pattern Recognition (CVPR)}, 2021.

\bibitem{cbdice}
Shi, P., Li, J., et al.\
``Centerline Boundary Dice Loss for Vascular Segmentation.''
\textit{Medical Image Computing and Computer-Assisted Intervention (MICCAI)}, 2024.

\bibitem{adamw}
Loshchilov, I., Hutter, F.\
``Decoupled Weight Decay Regularization.''
\textit{ICLR}, 2019.

\bibitem{gpt3}
Brown, T. B., Mann, B., Ryder, N., et al.\
``Language Models are Few-Shot Learners.''
\textit{Advances in Neural Information Processing Systems (NeurIPS)}, 2020.

\bibitem{bert}
Devlin, J., Chang, M.-W., Lee, K., Toutanova, K.\
``BERT: Pre-training of Deep Bidirectional Transformers for Language Understanding.''
\textit{NAACL-HLT}, 2019.

\bibitem{flashattn}
Dao, T., Fu, D. Y., Ermon, S., Rudra, A., R\'e, C.\
``FlashAttention: Fast and Memory-Efficient Exact Attention with IO-Awareness.''
\textit{NeurIPS}, 2022.

\bibitem{flashattn2}
Dao, T.\
``FlashAttention-2: Faster Attention with Better Parallelism and Work Partitioning.''
\textit{International Conference on Learning Representations (ICLR)}, 2024.

\bibitem{flashattn3}
Shah, J., Bikshandi, G., Zhang, Y., Thakkar, V., Ramani, P., Dao, T.\
``FlashAttention-3: Fast and Accurate Attention with Asynchrony and Low-Precision.''
\textit{NeurIPS}, 2024.

\bibitem{pytorch_amp}
PyTorch Documentation.\
\textit{Automatic Mixed Precision examples} (recommended bf16 pattern without
\texttt{GradScaler}). \\
\url{https://pytorch.org/docs/stable/notes/amp_examples.html}

\end{thebibliography}

\end{document}
